% Core written proof for checker_substitutive.py.  This section is independent
% of the earlier any-variable alternative and uses NEWEST witnesses only.
\section{Three instructions with no supplied introduction witness}
\label{sec:three-rule-core}

\subsection{Numerical notation and its exact reductions}

The numerical syntax is $0$, free and bound numerical variables, and
\[
                       R(n,a;i,u.h),
\]
where only $i,u$ are bound in $h$.  The defining computations are
\[
 R(0,a;i,u.h)=S(a),\qquad
 R(Sn,a;i,u.h)=h[n/i,R(n,a;i,u.h)/u],
 \qquad S(a):=R(0,a;i,u.0).
\]
Zero-counter bodies are canonically erased: $R(0,a;i,u.h)$ is represented
as $R(0,a;i,u.0)$.  Terms are full computation normal forms.  Every
substitution below is simultaneous, capture avoiding, and followed by
normalization, written $N$.  This is a metasyntactic convention; no
object-language universal evaluator or uniform normalization bound is
postulated.

The following three finite macros are useful:
\begin{align*}
 P(n)&=R(n,0;i,u.i),\\
 C(n,a,b)&=P(R(n,a;i,u.S(b))),\\
 J(n,a;i,u.h)&=P\bigl(R(n,a;j,v.S(h[j/i,P(v)/u]))\bigr).
\end{align*}
Here $j,v$ are fresh and substitution in the last display is simultaneous.
Equivalently its update is $(j,v)\mapsto S(h(j,P(v)))$.
All of the following are literal identities after normalization:
\begin{align*}
 P(Sn)&\equiv n,&P(0)&\equiv1,\\
 C(0,a,b)&\equiv a,&C(Sn,a,b)&\equiv b,\\
 J(0,a;i,u.h)&\equiv a,
 &J(Sn,a;i,u.h)&\equiv h[n/i,J(n,a;i,u.h)/u].
\end{align*}
In particular, $P$ is a left inverse of successor; its zero value is
\emph{one}, not zero.  When a conventional predecessor is wanted, use
$p(n)=C(n,0,P(n))$.  It has $p(0)\equiv0$ and $p(Sn)\equiv n$.
The scalar iterator of the source calculus embeds as
$I(n,a;u.h)=J(n,a;i,u.h)$ with the counter coordinate unused.

\subsection{Complete logical rules}

Formulas are equations and $D_{x,y}B$.  A clause is either $[\,;y.B]$
or $[A;y.B]$, with $y$ absent from its antecedent.  Every derivation starts
with the one fixed clause
\[
                        \mathcal R=[\,;y.y=y].
\]
An ordinary local assumption $A$ is the unconditional constant clause
$[\,;v.A]$.  The ordinary readings of $D$ and a conditional clause are
\[
 D_{x,y}B\ :=\ \exists x\bigl(B(x,0)\to\forall y B(x,y)\bigr),
 \qquad [A;y.B]\ :=\ A\to\forall yB.
\]

An environment orders its free numerical parameters.  Let $\ell$ be its
newest parameter, and let $\ell=0$ in an empty environment.  For a matrix
$B$ define $B_r(y)=N(B(r,y))$.  The three instruction families are:
\begin{enumerate}
\item \textbf{Use}: select a retained clause and a well-scoped term $t$;
prove its antecedent when it has one; conclude its normalized body at $t$.
\item \textbf{Introduce}: conclude $D_{x,y}B$ from
\[
 \Gamma,B_\ell(\eta)\vdash B_\ell(Sz),
 \qquad
 \eta=\begin{cases}
 z,&B_\ell(0)\text{ is literally an equation }r=r,\\
 0,&\text{otherwise},
 \end{cases}
 \tag{I}\label{eq:new-i}
\]
where $z$ is fresh.  The selected witness is the newest \emph{outer}
parameter; the new $z$ is available only in the premise.  No witness,
induction mode, or base proof is supplied in the instruction.
\item \textbf{Substitute and discharge}: prove a major $D_{x,y}B$ and
open its ordinary clause at a fresh parameter $b$:
\[
             \mathcal C_B(b)=[N(B(b,0));y.N(B(b,y))].
\]
In that same extended context prove $s=t$ and $N(F(s))$, for a supplied
one-place native formula template $F$.  Conclude $N(F(t))$, provided that
this conclusion contains no free $b$:
\[
 \frac{\Gamma\vdash D_{x,y}B\quad
       \Gamma,\mathcal C_B(b)\vdash s=t\quad
       \Gamma,\mathcal C_B(b)\vdash N(F(s))}
      {\Gamma\vdash N(F(t))}.
 \tag{E}\label{eq:new-e}
\]
The terms $s,t$ and the template may use $b$.  The final normalized
formula may not.  The template's placeholder is separate from all
numerical binders; substitution into it is capture avoiding.
\end{enumerate}
There are no other logical, equality, arithmetic, case, induction, or
structural rules.  The computations stated above belong to canonical
syntax.  The designated contradiction is $0=S0$.

\subsection{Soundness and the role of the selector}

Use is ordinary clause application.  Rule~\eqref{eq:new-e} is ordinary
existential elimination followed by equality substitution in its branch;
the condition on the final formula is exactly its eigenparameter condition.
This remains sound when a substitution target contains $b$.

For~\eqref{eq:new-i}, a nonreflexive zero row selects cases, and numerical
coverage combines the retained zero row with the supplied successor
derivation.  A reflexive zero row selects progress, but only for a literal
equation; contextual equation induction proves all its instances.  Its
zero instance is a reflexivity theorem.  In either case $\ell$ supplies
a legal existential witness.  Thus every rule is sound in ordinary
intuitionistic arithmetic with coverage, separation, and equation induction.
In particular the selector introduces no induction on quantified formulas.

At an allowed witness, a case premise can always be used.  If the selector
instead requests progress, its discarded zero-row assumption is literally
reflexive.  Supply that theorem in place of the old assumption and ignore
the new induction assumption.  In the constructive proof below this is
used as an operation on uniformly supplied branch builders, not as an
unproved parameter-weakening principle for arbitrary target proof trees.

\subsection{Preparatory constructions independent of equality transport}

Write $q_yA=D_{w,y}A$ with $w$ unused.  Its introduction uses an unused
witness, so changing which outer parameter is newest does not change its
matrix instances.  Its ordinary elimination follows from E by using
$s=t=0$ and the constant template whose value is the wanted conclusion.
The equality premise is Use of $\mathcal R$ at $0$.

Consequently the ordinary case-introduction and elimination rules for $q$
are constructible from I, E, and Use.  In a case introduction, if its zero
row is reflexive, use the preceding zero-row observation.  This gives
numerical cases for all native formulas, with branch proofs constructed
in their exact environments.  Neither arithmetic equality transport nor
equation induction is needed for these logical case constructions.

For completeness, formula gates and padding are likewise available before
transport.  Gate every numerical side of the bottom equation in $A$ by
$C(v,0,-)$, leaving its $D$ binders intact; call the result $G_vA$.
Then $G_{Sz}A\equiv A$, while $G_0A$ is a finite tower of $D$ binders
over $0=0$, and is uniformly provable by I and reflexivity.
For persistent premise builders, the padding $L(A)=q_yG_yA$ is
interderivable with $A$: one direction
ignores its local base and uses $A$, and the other uses the closure at
$1$ with the uniformly provable zero row.  Padding to the same binder
depth allows the usual componentwise selector of two formulas.  Thus
implication is $q_y\operatorname{sel}(y,A,B)$, up to this harmless
padding, with its ordinary introduction and elimination constructions.
This paragraph describes uniform proof-building operations; arbitrary
target cut or parameter weakening is only recovered after the main
equivalence theorem. Every introduction used here has an unused witness,
except introductions
of the all-zero truth towers, whose result is independent of that witness.

\subsection{An equality-free name for every numerical term}

Define
\[
 \operatorname{Name}(t)
   =D_{x,u}\bigl(C(u,0,x)=C(u,0,t)\bigr).
\]
Ordinary elimination at a fresh $b$ supplies $b=t$ by Use at $1$;
the needed antecedent is literally $0=0$.  No equality reasoning is
involved in this extraction.

Put
\[
 F(y)=D_{x,u}\bigl(
 C(y,0,C(u,0,x))=C(y,0,C(u,0,P(y)))\bigr).
\]
Its zero instance is a tower over $0=0$, so I at whichever outer
parameter is newest proves $F(0)$.  In the successor branch, $z$ is
the newest parameter.  Introduce the inner $D$ with the automatically
selected witness $z$.  Its zero row is $0=0$ and its successor row is
\[
                         z=P(Sz)\equiv z=z.
\]
Use of $\mathcal R$ proves this row, ignoring the local hypothesis.
Numerical cases therefore prove $F(S(t))$.  All computations are
literal: $P(S(t))\equiv t$ and $C(S(t),a,b)\equiv b$.  The resulting
formula is exactly $\operatorname{Name}(t)$.

This proof is reconstructed in each ambient environment.  Its important
features are that it uses no arbitrary introduction witness, no equality
substitution, and no predecessor theorem.  The shift by $S(t)$ is necessary;
the zero value $P(0)=1$ is erased by the outer gate.

\subsection{Translation from the diagonal native kernel}

Use the established native diagonal comparison calculus as source.
Its terms include the scalar iterator, it has ordinary equality transport,
closure use and witness elimination, and its $D$ introduction has an
arbitrary supplied witness.  Its progressive matrices have the literal
shape
\[
                   I(y,y;u.h)=I(y,y;v.k).
\]
All other introductions are by cases.  The repository proves this source
equivalent to the unrestricted native calculus, and hence to its ordinary
iPRA realization.  Map its scalar iterator to $J$ with the counter unused,
map its equations and binders compositionally, and normalize numerical
subterms.  A progressive diagonal matrix always has the literal normalized
zero row $0=0$, independently of its witness and its update bodies.

The translation is a relative induction on source derivations.  Its
induction statement allows an arbitrary target ambient environment and
context, with persistent builders for the source assumptions.  A builder
can be replayed after adjoining fresh parameters and clauses.  At the
root an assumption builder is hypothesis use.  Every replacement of a
local source assumption below uses only hypothesis use, reflexivity,
clause use, and transport through a dummy major
\[
                              D_{x,y}(0=0).
\]
The dummy major's introduction is independent of its selected witness,
so these builders remain persistent.  This invariant replaces the
otherwise invalid shortcut of weakening arbitrary target proofs after
the newest parameter has changed.

Computation axioms translate to reflexivity, because their two sides
have the same normal form.  Hypothesis use invokes its persistent builder.
Ordinary source witness elimination is an E step with a constant template
and the reflexive equality $0=0$.  Its local source assumption $Q_yB(b,y)$
is represented by the target clause $[B(b,0);y.B(b,y)]$.  Reify that clause
as the native formula $q_yB(b,y)$ by an unused-witness I instruction:
in its branch use the retained clause at $Sz$, supplying its zero antecedent
from the local case assumption, or from reflexivity if the selector chose
progress.  This reification is a persistent builder because its selected
witness is unused.  It supplies exactly the local-assumption realization
required by the induction hypothesis for the source elimination branch.

For a source $Q$ elimination, translate its major to $q_yB(y)$ and open
that formula first.  Only then invoke the induction hypothesis on the
strictly smaller source proof of $B(0)$, in the extended target context.
Use its resulting proof as the antecedent of the newly retained clause,
specialize at the requested term, and close E with a constant template
and $0=0$.  One must not translate the base premise outside the opening
and then silently weaken the resulting target proof tree into its branch.

For a source equality-transport node, first open the dummy major above.
Now apply the induction hypothesis to its two strictly smaller source
premises in that extended context, and finish E with their prescribed
formula template.  If the source uses the crossed one-term convention,
derive symmetry through the same dummy construction, or use the equivalent
full-transport source presentation.  It would be circular at this stage
to claim arbitrary target parameter weakening in order to move already
translated premise trees into the dummy branch; reapplying the induction
hypothesis is the correct construction.

For a source introduction at witness $t$, first prove
$\operatorname{Name}(t)$ and open it at a fresh $b$.  The branch receives
the persistent equality $b=t$, and $b$ is now the newest outer parameter.
The target I instruction therefore selects exactly $b$.

If the source introduction is progressive, its matrix is diagonal.
Its normalized zero row is $0=0$ even at $b$, so target I also selects
progress.  Transport the target assumption at $b$ to the source
assumption at $t$, use the induction hypothesis on the smaller source
branch, and transport its result back to $b$.

If the source introduction is by cases and target I chooses cases,
perform the same two transports on the zero-row assumption and successor
conclusion.  If target I chooses progress instead, its zero row at $b$
is reflexive.  Obtain that row by Use, transport it to the zero row at
$t$, use the induction hypothesis on the smaller source branch, and
transport back.  The extra target progress assumption is unused.
For the final transport back, open its dummy major before invoking the
induction hypothesis on the source branch; the recursive call is made in
that actual further-extended context. The builders for its assumptions
remain persistent there. Thus no completed target branch proof is moved
into a new parameter scope by an unproved weakening step.

In every case I produces the unchanged major formula $D_{x,y}B$, which
does not contain $b$.  Close its naming elimination with a constant
template and $0=0$.  Every recursive translation call is on a strictly
smaller source derivation; there is no appeal to the node being replaced.

\begin{theorem}[Consequence preservation]
The new three-instruction calculus proves exactly the normalized
translations of the native diagonal kernel, with the ordinary logical
and arithmetic meanings described above.  The transformation preserves
open assumptions through their stated realizations, and in particular
preserves closed theorems and the contradiction $0=S0$.
\end{theorem}

For target formulas using $R$ essentially, extend the comparison language
by that fixed primitive-recursive definition before applying the argument.
This is a definitional extension of the ordinary iPRA realization.  The
logical comparison and the diagonal argument are parametric in additional
numerical terms: full formula transport is already permitted, and the only
special numerical identities used by them are those of $0,S,I$.
The additional $R$ computation axioms again translate to reflexivity.
This accounts for the whole target language, not merely its scalar-iterator
image.  The written parametric argument is not a claim that the unchanged
Lean files have already formalized the enlarged syntax.

\subsection{Whole-domain interpretations and both numerical composites}

Let $T_R$ be ordinary intuitionistic first-order equality logic over the
staged $R$ descriptions, with their displayed computations, numerical
coverage, $(0=S0)\to\bot$, and contextual induction for literal equations.
The native logical comparison supplies its complete proof presentation.
The explicit falsity bridge in this ordinary realization is essential;
the equation $0=S0$ is the native contradiction, not by itself an ordinary
logical falsum in a positive singleton structure.

Interpret $R$ in iPRA by its fixed primitive-recursive equations.
Conversely, initial functions and composition translate directly, and a
fixed primitive recursion
\[
 f(0,\vec p)=g(\vec p),\qquad
 f(Sn,\vec p)=h(n,f(n,\vec p),\vec p)
\]
translates to
\[
                f^*(n,\vec p)=J(n,g^*(\vec p);i,u.h^*(i,u,\vec p)).
\]
There is no pairing construction and no universal evaluator here: the
counter coordinate is an explicit bound argument of each fixed update.
The ordinary predecessor $p$, selector $C$, addition and monus have their
usual equations.  The inherited equality-reflection argument from coverage,
separation and literal-equation induction then yields every quantifier-free
induction instance.  Successor injectivity follows immediately by applying
$P$; successor separation follows by the selector's zero/successor test.

For two fixed recursive definitions with identical initial values and
provably identical updates, literal-equation induction on equality of their
values proves uniqueness at every internal counter.  Apply this internally
while inducting externally on each finite description: a double-translated
description has the same equations as its original, and hence the same
values.  This argument works for both the $R$ and ordinary primitive-recursive
descriptions, with every passive parameter fixed during each induction.

Both interpretation domains are the entire numerical sort, defined by
$x=x$, and equality is actual equality.  The two numerical comparison
isomorphisms are therefore both the identity graph $x=y$.  The displayed
meanings of $D$, clauses, and the derived logical connectives give the two
formula round trips by external induction on formula syntax.  Preservation
of $0=S0$ and ordinary falsity gives effective proof translations of
contradictions in both directions and therefore equiconsistency.

These are internal uniqueness and formula-equivalence arguments, not an
inference from agreement in the standard natural numbers.  No claim of an
absolute minimal presentation or a new Lean verification is made by this
proof alone.
